\documentclass[11pt]{article}
\usepackage[utf8]{inputenc}
\usepackage[T1]{fontenc}
\usepackage[ngerman]{babel}
\usepackage{url}
\usepackage[margin=2.5cm]{geometry}

\title{Algorithm Engineering\\Woche 1, Aufgabe 1}
\author{Jan Vogt}
\date{}

\begin{document}

\maketitle

\section*{Moore's Law (Folie 10)}

Auf Folie 10 geht es um die Frage, warum man heute überhaupt parallel programmieren muss. Dort
steht "`The Free Lunch Is Over"' und darunter Moore's Law. In diesem Bericht gehe ich genauer auf
Moore's Law ein und erkläre, wie es mit den vielen Kernen in heutigen Prozessoren zusammenhängt.

\subsection*{Was ist Moore's Law?}

Gordon Moore, einer der Gründer von Intel, hat 1965 einen Artikel in der Zeitschrift Electronics
veröffentlicht. Darin hat er beobachtet, dass sich die Anzahl der Transistoren auf einem Chip
ungefähr jedes Jahr verdoppelt. Zehn Jahre später, also 1975, hat er seine Vorhersage auf eine
Verdopplung alle zwei Jahre angepasst. In dieser Form kennt man Moore's Law auch heute noch.

Eigentlich ist es kein richtiges Gesetz wie in der Physik, sondern eher eine Beobachtung, die
aber erstaunlich lange gestimmt hat. Der erste Mikroprozessor von Intel, der 4004 aus dem Jahr
1971, hatte gerade einmal etwa 2300 Transistoren, während große Chips heute mehrere zehn
Milliarden davon haben.

Oft liest man auch, dass sich die Leistung von Prozessoren alle 18 Monate verdoppelt. Das hat
Moore selbst so aber nie gesagt. Die 18 Monate stammen von einem Kollegen bei Intel, der
zusätzlich mit eingerechnet hat, dass die Transistoren mit der Zeit auch schneller werden.

\subsection*{Warum werden Programme nicht mehr von alleine schneller?}

Früher hat man von kleineren Transistoren automatisch profitiert, weil sie schneller schalten
und gleichzeitig weniger Strom brauchen. Dieser Zusammenhang wird nach Robert Dennard (1974) als
Dennard Scaling bezeichnet. Dadurch konnte man die Taktfrequenz immer weiter erhöhen, ohne dass
der Chip zu heiß wurde. Programme sind also mit jeder neuen Prozessorgeneration schneller
geworden, ohne dass man als Programmierer etwas dafür tun musste. Genau das ist mit dem
"`Free Lunch"' gemeint.

Um das Jahr 2005 hat das Dennard Scaling dann aufgehört zu funktionieren. Bei sehr kleinen
Transistoren fließt nämlich auch dann Strom, wenn sie eigentlich ausgeschaltet sind
(Leckströme), und die Chips wurden dadurch zu heiß. Deshalb ist die Taktfrequenz seitdem kaum
noch gestiegen. Der Pentium 4 hatte 2004 schon 3,8 GHz und normale Prozessoren liegen heute nicht
viel höher.

Moore's Law gilt aber weiterhin, es kommen also immer noch mehr Transistoren dazu. Die Frage war
nur, was man mit den ganzen zusätzlichen Transistoren anfangen soll, wenn man den Takt nicht mehr
erhöhen kann.

\subsection*{Die Formel auf der Folie}

\[
  \mathrm{power} = \mathrm{performance}^{1.74}
\]
Die Zahl kommt aus Messungen, die Intel an den eigenen Prozessoren gemacht hat (Grochowski und
Annavaram). An der Formel sieht man ganz gut, warum es sich nicht lohnt, einen einzelnen Kern
immer schneller zu machen. Soll ein Kern doppelt so schnell sein, braucht er
$2^{1.74} \approx 3{,}3$ mal so viel Strom. Nimmt man stattdessen zwei Kerne mit der alten
Geschwindigkeit, verdoppelt sich der Stromverbrauch nur ungefähr und man hat trotzdem die
doppelte Leistung. Das funktioniert allerdings nur, wenn sich das Programm auch parallelisieren
lässt.

Aus diesem Grund bauen die Hersteller seit etwa 2005 Prozessoren mit mehreren Kernen. Der Rechner
im PC-Pool hat zum Beispiel 20 Kerne (8 P-Cores und 12 E-Cores).

Herb Sutter hat dazu 2005 den Artikel "`The Free Lunch Is Over"' geschrieben. Seine Aussage ist,
dass Programme nicht mehr von alleine schneller werden und man sie selbst so schreiben muss, dass
sie mehrere Kerne nutzen, wenn man mehr Leistung haben will. Genau dafür lernen wir in der
Vorlesung OpenMP.

\subsection*{Gilt Moore's Law heute noch?}

Darüber wird ziemlich viel diskutiert. Die Verdopplung dauert inzwischen eher länger als zwei
Jahre und neue Fertigungsprozesse werden immer teurer. Außerdem sind Namen wie "`3 nm"' eher
Marketing und haben mit einer echten Größe auf dem Chip nicht mehr viel zu tun. Der Chef von
Nvidia hat 2022 zum Beispiel gesagt, dass Moore's Law tot ist, während Intel das anders sieht.

\subsection*{Fazit}

Wichtig ist, dass Moore's Law eigentlich nur etwas über die Anzahl der Transistoren aussagt
und nicht über die Geschwindigkeit. Mehr Transistoren gibt es zwar weiterhin, aber heute bekommt
man sie vor allem in Form von mehr Kernen. Damit ein Programm davon auch etwas hat, muss man es
selbst parallelisieren.

\end{document}
